\section{The ordinary extension by finite linear orders}
\label{sec:complete-new-model}

This section reports the second model studied in the workspace.  Starting
with the OpenAI symmetric model $W$, we force a linear extension of its
native partial order and obtain $N=W[G]$.  The established conclusions
include
\[
 N\models\mathsf{ZF}+\mathsf{AC}_{\mathrm{WO}}+\mathsf{DC}
       +\neg\mathsf{AC}+\neg\mathsf{DC}_{\omega_1},
 \qquad \mathsf{SVC}(A)+\mathsf{SVC}(\mathbb R_N),
\]
where $\mathbb R_N=\mathcal P(\omega)^N$.  All old cardinal comparisons
survive, but the new real set is not well-orderable.  Whether $N$ satisfies
\emph{global} PP, NDS, FB, or BPI remains unresolved.  In particular, this
section does not present a model of BPI together with failure of Choice.

Here $X\approx Y$ means equipotence, $X\preccurlyeq Y$ means the existence
of an injection, and ``old'' means belonging to $W$, rather than merely
to the earlier ZFC construction ground $V$.  A result described as
audited has a written mathematical proof and a separate workspace
check; this is neither external peer review nor a formal verification
of the whole forcing construction.  The proofs below are summaries of
the saved arguments.  The audit register at the end supplies their
exact filenames.

\subsection{Construction and the chain conditions}

The starting model has a non-well-orderable set $A$ and a ground monoid
$T$ of well-orderable cardinality $\mu$.  There is an internal free right
action $(a,t)\mapsto a\cdot t$ of $T$ on $A$.  ``Free'' here means that
$t\mapsto a\cdot t$ is injective for each fixed $a$; it does not mean
that right multiplication $a\mapsto a\cdot t$ is injective.  Define
\[
 a\leq_T b\quad\Longleftrightarrow\quad
       (\exists t\in T)\ b=a\cdot t.
\]
This is a partial order.  Its principal upset $aT$ is called a
\emph{chart}; it is an old well-orderable set of cardinality $\mu$.
Every old well-orderable subset of $A$ lies in one chart.  The results
used from $W$ include
\[
 A^2\approx A,\qquad A\times\mu\approx A,\qquad
 (A^\omega)^W\approx A,\qquad h^W(A)=\mu^+.
\]
Every old set has a normal form $(A\times\alpha)\sqcup\beta$, with
ordinal coefficients.  Here $h(X)$ is the least ordinal not injecting
into $X$; $w(A)=\mu$ means that $\mu$ is the largest well-orderable
cardinal represented by a subset of $A$.

Let $\mathbb P_{\mathrm{ext}}$ consist of finite strict linear orders on
finite subsets of $A$ which extend the strict part $<_T$ of the native order on their domains.
A stronger condition extends the domain and the finite order.  The
generic union $<_L$ is a strict linear order on all of $A$ extending $<_T$; we also write $L$ for this order.
Set
\[
              N=W[G],\qquad G\subseteq\mathbb P_{\mathrm{ext}}.
\]
This is ordinary set forcing over the choiceless model $W$.  Its
construction and its preservation properties require separate proofs.

\paragraph{Audited chain-condition theorem.}
Every antichain in $\mathbb P_{\mathrm{ext}}$ is countable, and every
predense subset has a countable predense subset.  We call these
properties $\mathrm{CCC}_2$ and $\mathrm{CCC}_3$, respectively.  The
second is not inferred from the first in arbitrary ZF models.

The local proof uses the source monoid's height $\rho:T\to(0,1)$,
positive costs $\epsilon_t$, and conditional joins.  Two conditions
$(x_0<\cdots<x_{n-1})$ and $(y_0<\cdots<y_{m-1})$ are incompatible
exactly when a reversed rectangle occurs:
\[
 i<j,\ l<k,\qquad x_j\leq_T y_l,
                         \quad y_k\leq_T x_i.
\]
After bounding the finite lengths and bounding the costs away from zero,
finite Ramsey theory turns a sufficiently large antichain into a
triangular pattern.  Successive finite joins would then increase
$\rho$ by more than a fixed positive amount at every step, contradicting
$0<\rho<1$.  This proves a local $\sigma$-bounded chain condition on a
chart.  An alleged $\omega_1$-indexed global antichain localizes to one
chart.  In $W$ every uncountable set contains $\omega_1$, so this excludes
all uncountable antichains.

Each chart subforcing is regular in $\mathbb P_{\mathrm{ext}}$.
A finite condition has a finite reduction to the chart: its constraints
are projected to the least native successors in that chart, using
conditional joins.  The stronger $\mathrm{CCC}_3$ proof then uses the
actual source construction's petal-separation lemma.  If a predense set
has supporting diagram $p$, choose a larger diagram $q$ with $q_e=p$.
Its part inside the chart of $q$ is globally predense.  A greedy
well-ordered construction in that chart produces a countable predense
subfamily.  The source-specific petal and gate calculations are material
hypotheses of this argument [N1].

\subsection{Preservation and the remaining failures of Choice}

\paragraph{Audited preservation theorem.}
The extension preserves all ordinal cardinals and cofinalities, the
Hartogs number of every old set, and the non-well-orderability of every
old non-well-orderable set.  More precisely, every new well-orderable
subset of an old set $X$ is contained in an \emph{old} well-orderable
subset of $X$ [N2].

For the proof let $\mathbb B=\operatorname{RO}^W(\mathbb P_{\mathrm{ext}})$.
Its antichains are countable: an uncountable Boolean antichain would
contain an $\omega_1$-indexed subfamily, and old
$\mathsf{AC}_{\mathrm{WO}}$ would select incompatible forcing conditions
below its members.  If a name is forced below $b_0$ to be a function
$f:X\to Y$ between old sets, put
\[
 g(x)=\{y\in Y:b_0\wedge\Vert\dot f(\check x)=\check y\Vert>0\}.
\]
For fixed $x$ the nonzero Boolean values are disjoint.  Thus $g(x)$ is
countable in $W$, and $g$ is one canonically defined old function with
$f(x)\in g(x)$.  For an ordinal domain, old
$\mathsf{AC}_{\mathrm{WO}}$ well-orders the union of these possible-value
sets.  This gives the asserted covering property, from which the
cardinal, cofinality and Hartogs conclusions follow.  In particular
\[
                w^N(A)=\mu,\qquad h^N(A)=\mu^+,
                \qquad N\models\neg\mathsf{AC}.
\]
The complete-Boolean-algebra covering argument is related to the
choiceless chain-condition analysis of Karagila--Schweber; the saved
proof checks the extra hypotheses needed here, instead of importing
ZFC preservation wholesale.\footnote{A.~Karagila and N.~Schweber,
``Choiceless chain conditions,'' \emph{European Journal of Mathematics}
8, Supplement 2 (2022), S393--S410,
\url{https://doi.org/10.1007/s40879-022-00564-2}.}

The stronger chain condition separately preserves
$\mathsf{AC}_{\mathrm{WO}}$.  For each member of a new ordinal-indexed
family, use a countable predense set of conditions carrying witness
names.  Collection first bounds those names in one old set.  Old
ordinal-indexed choice selects the countable lists and witnesses;
in the generic extension, the first successful condition supplies the
chosen element.  This uses local witness density, not the choiceless
maximum principle.  Consequently $N$ also satisfies ordinary DC,
by the established implication $\mathsf{AC}_{\mathrm{WO}}\Rightarrow
\mathsf{DC}$ [N2].

Every new ordinal sequence of $A$-points is contained in one old chart.
It therefore has a strict common native predecessor.  A descending
$\omega_1$-sequence would produce, in that chart, a strictly decreasing
$\omega_1$-sequence of real heights, which is impossible.  The resulting
countably closed tree of descending native sequences witnesses
$\neg\mathsf{DC}_{\omega_1}$.  Also $L$ has no well-orderable coinitial
subset.  Its reverse witnesses failure of the Cofinality Principle,
the assertion that every linear order has a well-orderable cofinal
subset.  These statements concern new subsets as well as old ones.

The old Boolean algebra $\mathbb B$ has cardinal $A$, is complete,
atomless and ccc in $W$, and has no ultrafilter there: an ultrafilter
would evaluate its comparison values to a forbidden native linear
extension.  The generic extension supplies an ultrafilter on this
\emph{old} algebra.  Its old completeness need not extend to all new
families of Boolean elements.  Repairing this particular obstruction
does not establish BPI for all Boolean algebras of $N$ [N3].

\subsection{The new real set and the unresolved evaluation quotient}

\paragraph{Audited real-cardinal theorem.}
In $N$ there are maps
\[
            A\lhook\joinrel\longrightarrow\mathbb R_N,
       \qquad A\twoheadrightarrow\mathbb R_N.
\]
The real set is not well-orderable and satisfies
\[
 w(\mathbb R_N)=\mu,\qquad h(\mathbb R_N)=\mu^+,
 \qquad \ell(\mathbb R_N)=\mu^+,
 \qquad (A^\omega)^N\approx\mathbb R_N.
\]
For the original parameters $\mu=\aleph_2$, this gives the non-well-orderable
real set Hartogs number $\aleph_3$; parameter variants use their preserved
$\mu$.  Here $\ell$ is the least positive ordinal not obtainable as a
surjective image; the harmless empty-set convention does not affect
the displayed infinite value [N4].

For the injection, fix selected probes $m_n$ and encode the order
pattern of $\{a\cdot m_n:n<\omega\}$ as a real $r_a$.
The probe geometry implies that the sets at two distinct centers
intersect in at most one point.  A finite-order swap can change a
comparison at one center while keeping the corresponding comparison
at the other unchanged.  These separating extensions are dense, so
$a\mapsto r_a$ is injective.

For the surjection, $\mathrm{CCC}_3$ expresses every old Boolean value
as a countable join of forcing conditions.  Since
$|\mathbb P_{\mathrm{ext}}|=|A|$ and $(A^\omega)^W\approx A$, one old
$A$-sized set codes all old sequences of Boolean values.  Evaluation
under the generic ultrafilter covers every new subset of $\omega$.
This produces a surjection, not an injection in the opposite direction.

If $\mu^+$ injected into $\mathbb R_N$, new
$\mathsf{AC}_{\mathrm{WO}}$ would choose preimages under
$A\twoheadrightarrow\mathbb R_N$, contradicting the preserved Hartogs
number of $A$.  This proves the upper threshold; the profile injection
gives the lower bound.  Finally,
$\mathbb R_N\preccurlyeq A^\omega\preccurlyeq
\mathbb R_N^\omega\approx\mathbb R_N$ proves the countable-power
identity.  The left-hand power here is computed in $N$, not in $W$.

The central unresolved comparison is
\[
                         \mathbb R_N\preccurlyeq A\ ?
\]
An injection and a surjection from $A$ to the same target do not
settle this in ZF.  Least choices of ground code \emph{types} do not
choose centers in the remaining non-well-orderable fibers.  Even
$\mathbb R_N\approx A$, if eventually proved, would settle one
quotient instance and would not alone prove global PP.

\subsection{Local real fields and exact localization}

Put $Q=\mathbb P_{\mathrm{ext}}(T)^V$.  For each $a\in A$, the chart
map $t\mapsto a\cdot t$ pulls the generic order back to a $Q$-generic
filter $K_a$.  Chart regularity gives genericity over $W$, hence over
$V$; it does not justify genericity over the external Cohen universe
$V[F]$.  Fix in $V$ a $\mu$-enumerated set $\Gamma$ of all countable
$Q$-names coding reals, and define internally
\[
 R_a=\{\tau^{K_a}:\tau\in\Gamma\}.
\]
Externally this is $\mathbb R_N\cap V[K_a]$.  The internal definition
uses a set of names, not a predicate for the ground universe.

\paragraph{Audited local-field theorem.}
For every $a$, $|R_a|=\mu$.  Every new well-orderable subset of
$\mathbb R_N$ is contained in some $R_a$.  Every new subset of an
ordinal belongs to $W[K_a]$ for some $a$.  Moreover
\[
 R_a\subseteq R_b\quad\Longleftrightarrow\quad a\in bT,
 \qquad
 r_a\in R_b\quad\Longleftrightarrow\quad a\in bT.
\]
Indeed the stronger exclusion $r_a\notin W[K_b]$ holds whenever
$a\notin bT$ [N5].

The size bound comes from the ground enumeration of $\Gamma$; the
$\mu$ profiles $r_{a\cdot t}$ give equality.  For localization, take
the Boolean bits of an ordinal-indexed real family and choose their
countable forcing codes using old ordinal-indexed choice.  All finite
point domains form one old well-orderable subset of $A$, hence lie
in one chart.  Each individual real then has a countable sequence of
ground $Q$-conditions.  The original model adds no countable sequences
of ground elements, so that code belongs to $V$ and evaluates in
$V[K_a]$.  The full uncountable array need only belong to $W$; this
explains the weaker $W[K_a]$ conclusion for arbitrary ordinal subsets.

The inclusion-reflection proof uses a quotient-forcing argument.  If
$a\notin bT$, infinitely many of the probes at $a$ have identical
native comparison profiles to the entire chart $bT$.  Below any finite
quotient condition, two remain interchangeable while all chart-order
information is fixed.  Both orientations can still be forced, so the
corresponding bit of $r_a$ is new over $W[K_b]$.

Two distinctions are essential.  A well-orderable real family is
\emph{contained in} $R_a$; the family itself need not belong to
$V[K_a]$.  Also proper field inclusions do not give strict cardinal
inequalities: all the $R_a$ have the same cardinal $\mu$.  Every
well-orderable family of these fields is contained in one strictly
larger field, by taking a strict native predecessor of its centers.
Thus there is no well-orderable cofinal subfamily under inclusion.

\subsection{All comparisons between old sets are preserved}

For any old $X,Y$, the existence in $N$ of an injection, bijection, or
surjection between them is equivalent to its existence in $W$.
Every new surjection between old sets has the PP conclusion witnessed
by an \emph{old} injection in the reverse direction.  Moreover every
new set-family consisting of old sets has a finite coinitial cardinal
basis [N6].

Here is the comparison invariant.  An old well-orderable set of
cardinality $\nu$ has record $(0,\nu)$.  Normalize a non-well-orderable
old set as
\[
 X\approx(A\times\alpha)\sqcup\beta,\qquad
 \alpha\geq\mu,\quad \beta=0\text{ or }\beta>\alpha,
\]
with ordinal cardinals as coefficients, and give it record
$(\alpha,\varepsilon)$, where $\varepsilon=\max(\alpha,\beta)$.
These records are unique and injection comparison is componentwise,
both in $W$ and in $N$.

For example, an injection from record $(\alpha,\varepsilon)$ into
$(\gamma,\zeta)$ gives $\varepsilon\leq\zeta$ by Hartogs preservation.
Each of the $\alpha$ source copies of $A$ must meet the target's
$A$-part; otherwise that copy would become well-orderable.  New
ordinal-indexed choice selects one point from each such disjoint
intersection, giving $\alpha\leq\gamma$.  Conversely the two
inequalities already construct an old injection.  A similar
preimage selection proves preservation of surjections.

Every set of ordinal pairs has a finite coinitial subset under the
componentwise order: repeatedly choose a least first coordinate and
then a least second coordinate not already dominated.  Continuing
infinitely would strictly decrease the second coordinate infinitely.
Applying this to the unique old records proves the relative FB claim.
For a new family of old sets, first bound their ranks inside one old
rank segment, so its old record assignment is a set function.

Thus an infinite new cardinal antichain or descending chain cannot
consist solely of old cardinalities.  This is not global FB or NDS
for all new sets.  If $\mathbb R_N\not\approx A$, its cardinality is
new: $A$ is the only old non-well-orderable cardinality with
$w=\mu$.

\subsection{Explicit nonsplitting and additional choice bounds}

For every $t\neq1$, the old surjection $f_t:A\to A$,
$f_t(x)=x\cdot t$, remains without a right inverse.  More strongly,
there is no assignment of a nonempty countable subset of each fiber,
even in $N$ [N7].  The old proof says that any countable-valued
assignment on $A$ has some value contained in the forward chart of
its argument.  A nonempty value consisting of predecessors would
then contradict conicality of $T$.  For a new assignment,
$\mathrm{CCC}_3$ supplies a canonical old countable cover of every
possible value, reducing the contradiction to that old statement.

Transport $f_t$ along $a\mapsto r_a$ and extend it by the identity
outside that image.  This gives a surjection
$g_t:\mathbb R_N\to\mathbb R_N$ without a right inverse, or even a
countable multisection.  Consequently choice fails for real-indexed
families of nonempty sets of reals.  This is not a failure of PP:
PP does not split each specified surjection, and these two sets are
already equipotent.

Ordinary forcing preserves $\mathsf{SVC}(A)$.  The profile injection
gives a retraction $\mathbb R_N\twoheadrightarrow A$, hence
$\mathsf{SVC}(\mathbb R_N)$.  Coding
$\mathbb R_N\times\eta$ into a power set of an ordinal yields
\[
             \mathsf{KWP}^{*}_1,\qquad \mathsf{KWP}_2
             \quad\hbox{in }N.
\]
Here $\mathsf{KWP}_n$ means that every set injects into an $n$-fold
power set of an ordinal; the starred version uses partial surjections.
Inverse-image fibers give $\mathsf{KWP}^{*}_1\Rightarrow
\mathsf{KWP}_2$.  These are the established bounds, not
$\mathsf{KWP}_1$, the Ordering Principle, or global finite-choice
preservation.\footnote{A.~Karagila and J.~Schilhan,
``Intermediate models and Kinna--Wagner Principles,''
\emph{Proceedings of the American Mathematical Society} 154 (2026),
393--403, \url{https://doi.org/10.1090/proc/17425}; see Section 3.2
and Proposition 4.1 of \url{https://arxiv.org/html/2409.07352v2}.}

There are also definite negative continuum-function conclusions.
Write ICF for
$\mathcal P(X)\approx\mathcal P(Y)\Rightarrow X\approx Y$,
and HCF for reflection of injection comparison by $\mathcal P$.
ZF proves $\mathsf{ICF}+\mathsf{KWP}_{n+1}\Rightarrow
\mathsf{KWP}_n$: apply the latter to $\mathcal P(X)$, obtaining
$\mathcal P(X)\preccurlyeq\mathcal P(B)$ with
$B=\mathcal P^n(\theta)$ and $\theta$ infinite.  Since $B$ absorbs
two copies,
$\mathcal P(X\sqcup B)\approx\mathcal P(B)$, so ICF gives
$X\preccurlyeq B$.  Finite descent reaches AC.  Thus the proved
$\mathsf{KWP}_2+\neg\mathsf{AC}$ package implies
$\neg\mathsf{ICF}+\neg\mathsf{HCF}$ in $N$ [N8].

\subsection{Initial segments, final segments, and open intervals}

For $a<_L b$ put
\[
 I_a=\{x:x<_L a\},\qquad J_a=\{x:a<_L x\},\qquad
 (a,b)_L=\{x:a<_L x<_L b\}.
\]
The final audited result is
\[
                  I_a\approx J_a\approx(a,b)_L\approx A.
\]
The earlier possibility that a final segment might be a proper
square root of $A$ was resolved by an explicit injection, not by
cancelling the established identity $J_a^2\approx A$ [N9].

An initial segment contains the old $A$-sized fiber
$\{x:x\cdot t=a\}$ for any $t\neq1$.  For an interval, choose selected
probes $m_n$ with $m_ne=e$, pairwise cone intersection $eT$, and only
identity common left divisors.  Fix $d$ with $d\cdot e=b$, and set
$F=\{x:x\cdot m_0=d\}\approx A$.  For each $x\in F$ the ray
\[
                    S_x=\{x\cdot m_n:n\geq1\}
\]
lies strictly below $b$.  These rays are disjoint for different $x$:
a shared ray point together with the shared anchor $d$ forces the
common-chart join quotients to divide both $m_0$ and a later $m_n$,
so both quotients are $1$.

In a finite condition, a point below $b$ can dominate at most one
point of $S_x$; dominating two would place it natively above $b$.
A remaining ray point can therefore be inserted immediately before
$b$.  Genericity makes $S_x$ cofinal below $b$.  Its least indexed
point above $a$ injects $F$ into $(a,b)_L$, proving the theorem.
Taking $b=a\cdot e$ also proves the final-segment case.

The whole $L$ has countable cofinality.  A fixed countable sequence
of fresh merger letters with disjoint cones gives, above any center,
an old countable cofinal ray.  Every old non-well-orderable subset of
$A$ is nevertheless $L$-coinitial: a point outside a chart covering
a finite condition can be placed first.  Thus no final segment or
bounded interval contains an old non-well-orderable subset.  The
$A$-sized copies just constructed are new.

All these segments have finite positive powers of cardinal $A$ and
new countable power of cardinal $\mathbb R_N$.  They therefore yield
no strict cardinal descent among themselves.  A further recorded
refinement says that each final segment has an old countable
order-dense subset, using the common-divisor and disjoint-cone
properties of the merger rays.  That refinement has a written proof
in [N9], but its separate independent audit is not part of the saved
segment audit; it is not used in any theorem above.

\subsection{Boundaries of the current classification}

The unresolved global questions in $N$ are PP, the strong quotient
property $Q(A)$, NDS, FB, the assertion that every cardinal antichain
is finite, and BPI.
Also unresolved are $\mathbb R_N\preccurlyeq A$,
$\mathsf{KWP}_1$, the Ordering Principle, and Choice for arbitrary
families of finite sets.  The original model's countable-member choice
consequence has not been transferred to $N$.  Adding a native linear
extension does not establish the Order Extension Principle for every
partial order, and neither it nor the local fields settles BPI.

Two bounded deductions sharpen the outstanding work.  First, a
hypothetical map $j:\mathbb R_N\to A$ has a support root $d$ such
that $j[R_a]\subseteq aT$ whenever $a\leq_T d$.  This produces a
coherence requirement among local $\mu$-sized injections, rather than
a cardinal contradiction.  Second, every directed ideal of $T$ in
$N$ has a canonically chosen countable cofinal sequence: select least
points in the old $T$-wellorder with heights approaching the supremum,
then take finite joins.  The positive-cost inequality rules out an
omitted divisor.  Hence all these ideals inject into
$\mu^\omega\approx\mathbb R_N$ [N10].  These are saved local
deductions; neither supplies the missing global coherent selector.

The original quotient proof explicitly requires support descent and
coherent matching across charts.  New availability ideals need not
be ground objects, and arbitrary ordered three-chart amalgams can
cycle.  A bound on local names or ideals by $\mathbb R_N$ is therefore
not a proof of an injective presentation theorem.  Likewise the
Truss-model comparison in Ryan-Smith's work shows why
$\mathsf{AC}_{\mathrm{WO}}+\mathsf{SVC}(\mathbb R)$ alone must not be
treated as PP.\footnote{C.~Ryan-Smith, ``Local reflections of choice,''
\emph{Acta Mathematica Hungarica} 176 (2025), 244--257,
\url{https://doi.org/10.1007/s10474-025-01533-3}.
The relevant comparison is Proposition 3.19 in the published version,
numbered 3.18 in the author-hosted preprint.  It concerns a different
model; it is not a negative PP theorem about $N$.}

\subsection*{Saved proof and status register for this section}

All filenames below are in the public archive directory
\path{evidence/original_model/}, mapped from the workspace record
\path{outputs/BPI_RPP_five_hour_research/}.  The source references are
the pinned OpenAI construction in \path{work/openai_math/}: in
particular \path{forcing.tex} (diagram amalgamation, local decision,
ordinal choice and no countable ground sequences), \path{monoid.tex}
(probe gates, merger normal forms and the construction schedule), and
\path{quotients.tex} (availability ideals and coherent quotient
matching).  The external construction is credited to
\url{https://github.com/openai/math}; the extensions and deductions
reported here must be distinguished from statements made in that
manuscript.

\begingroup\small
\noindent\textbf{[N1] Chain conditions and chart reduction, audited.}
\path{FINITE_LINEAR_EXTENSION_FORCING_CCC.txt};
\path{FINITE_LINEAR_EXTENSION_CCC_INDEPENDENT_AUDIT.txt};
\path{FINITE_ORDER_FORCING_CCC3_CHART_REDUCTION.txt};
\path{FINITE_ORDER_FORCING_CCC3_INDEPENDENT_AUDIT.txt}.

\noindent\textbf{[N2] Preservation, audited.}
\path{BOOLEAN_COMPLETION_COVERING_AND_HARTOGS_PRESERVATION_AUDIT.txt};
\path{CCC3_PRESERVES_AC_WO_CONDITIONAL_AUDIT.txt};
\path{AC_WO_DC_PRIMARY_SOURCE_RECONCILIATION.txt}.
The ``conditional'' filename records an earlier checkpoint: its
$\mathrm{CCC}_3$ hypothesis was subsequently verified in [N1].

\noindent\textbf{[N3] The old complete Boolean algebra, audited.}
\path{COMPLETE_CCC_BOOLEAN_ALGEBRA_EXACT_CARDINAL_A_AUDIT.txt};
\path{BOOLEAN_ALGEBRA_ROOT_AUDIT.txt}.

\noindent\textbf{[N4] Real maps and thresholds, audited.}
\path{GENERIC_ORDER_PROFILES_INJECT_A_INTO_REALS.txt};
\path{GENERIC_REAL_PROFILE_INDEPENDENT_AUDIT_AND_QUOTIENT_TEST.txt};
\path{BOOLEAN_CODES_FOR_NEW_REALS_ROOT_DRAFT.txt};
\path{BOOLEAN_CODES_REAL_QUOTIENT_INDEPENDENT_AUDIT.txt}.
The retained ``draft'' filename is accompanied by the separate audit.

\noindent\textbf{[N5] Local fields and localization, audited.}
\path{LOCAL_REAL_FIELDS_REFLECT_NATIVE_ORDER.txt};
\path{LOCAL_REAL_FIELDS_INDEPENDENT_AUDIT.txt};
\path{LOCAL_ZFC_REAL_COVERS_AND_ORDINAL_LOCALIZATION_AUDIT.txt}.

\noindent\textbf{[N6] Old comparisons and relative finite bases.}
\path{OLD_CARDINAL_COMPARISONS_AND_RELATIVE_FINITE_BASIS.txt};
\path{OLD_SURJECTIVE_COMPARISONS_AND_RELATIVE_PP.txt}.
These are mathematical deductions from the audited preservation and
normal-form results; the surjective comparison has a recorded second
reconstruction.

\noindent\textbf{[N7] Nonsplitting and multisections, audited.}
\path{COUNTABLE_MULTISECTIONS_AND_NONSPLITTING_PRESERVATION.txt};
\path{COUNTABLE_MULTISELECTOR_INDEPENDENT_AUDIT.txt}.

\noindent\textbf{[N8] Choice bounds and continuum-function failure.}
\path{NEW_MODEL_KWP_OP_FINITE_CHOICE_BOUNDARY.txt};
\path{FINITE_KWP_ICF_COLLAPSE_INDEPENDENT_AUDIT.txt}.

\noindent\textbf{[N9] Resolved segment cardinalities, audited.}
\path{ANCHORED_RAYS_RESOLVE_TAIL_AND_INTERVAL_CARDINALS.txt};
\path{GENERIC_ORDER_SEGMENTS_INDEPENDENT_AUDIT.txt}.
The preliminary file
\path{GENERIC_LINEAR_ORDER_SEGMENTS_AND_SQUARE_ROOTS.txt}
now carries a supersession notice.  Its open tail alternative is not
a current unresolved question.  The separability refinement in the
later note remains separately flagged above.

\noindent\textbf{[N10] Bounded further deductions and remaining gaps.}
\path{REAL_TO_LABEL_INJECTION_CHART_LOCALITY_AUDIT.txt};
\path{NEW_AVAILABILITY_IDEALS_REAL_CODE_BOUND.txt};
\path{TRUSS_QUOTIENT_COMPARISON_AFTER_ORDER_FORCING.txt}.
These do not assert a completed new quotient classification.
\endgroup
